\chapter{Pensieve: Neural Adaptive Video Streaming}
\paperinfo{H. Mao, R. Netravali, M. Alizadeh (MIT CSAIL)}{ACM SIGCOMM 2017}

\begin{keybox}[Pensieve in one paragraph]
A video player must pick a \key{bitrate for every chunk} of video (an \key{ABR}, adaptive bitrate, algorithm). Existing ABR algorithms use \key{fixed rules} built on simplified models (``predict throughput'', ``keep the buffer at a level'', or MPC's optimisation over a short horizon with a throughput forecast), so they cannot be best for every network and every QoE goal. \key{Pensieve} instead \key{learns} the ABR policy with \key{deep reinforcement learning}: a \key{neural network} maps raw observations (past chunk throughputs and download times, next chunk sizes, buffer level, chunks left, last bitrate) to a probability for each bitrate, and is trained with \key{A3C} (actor-critic) so that the cumulative \key{QoE reward} is maximised. Training runs in a fast \key{chunk-level simulator} over many network traces (100 hours of video in 10 minutes); \key{masking} lets one model serve videos with different bitrate ladders; at run time the model sits on an \key{ABR server} that clients query for each chunk.
\end{keybox}

\section{Background: HTTP adaptive streaming (DASH)}
\paperfig[0.6]{ps_fig1_dash}{HTTP adaptive streaming: the player gets a manifest, then requests chunks from a CDN one by one; the ABR controller chooses each chunk's quality from throughput estimates and buffer occupancy.}
\begin{itemize}
\item Video is delivered over \key{HTTP} (standardised as \key{DASH}): providers reuse existing \key{CDNs}, keep \key{stateless} back-ends, and work with browsers and apps.
\item A video is split into \key{chunks} of a few seconds. Each chunk is encoded at several \key{discrete bitrates}; higher bitrate means higher quality and a larger chunk. Chunks are \key{aligned across bitrates}, so the player can switch quality at any chunk boundary without redundant bits or skipped video.
\item Session flow: the player sends an \key{authentication token} to the provider; gets back a \key{manifest} listing the CDN and the available bitrates; then requests chunks one by one, choosing each chunk's bitrate with the ABR algorithm.
\item A chunk can only be \key{played after it is completely downloaded}. Downloaded video waits in the \key{playback buffer}; if the buffer empties, playback stalls (\key{rebuffering}).
\end{itemize}

\subsection{Why ABR is hard}
\begin{enumerate}
\item \key{Variable networks}: throughput fluctuates and differs across environments. On cellular links throughput prediction is often wrong, so stable signals such as \key{buffer occupancy} should get more weight there.
\item \key{Conflicting QoE goals}: high bitrate, little rebuffering, and smoothness (few quality switches) pull against each other, and users weigh them differently.
\item \key{Cascading effects}: a high bitrate now can drain the buffer and force low bitrates (or stalls) later; with a smoothness goal, today's choice also biases tomorrow's.
\item \key{Coarse decisions}: only a few bitrates exist. If throughput is just below one level but well above the next, the algorithm must gamble between quality and rebuffering risk.
\end{enumerate}

\subsection{Existing ABR algorithms}
\begin{center}
\begin{tabularx}{\linewidth}{@{}lXX@{}}
\toprule
\textbf{Family} & \textbf{Rule} & \textbf{Weakness} \\
\midrule
\key{Rate-based} (e.g.\ FESTIVE) & Predict throughput (e.g.\ \key{harmonic mean of the last 5 chunks}); pick the highest bitrate below it & Throughput estimates over HTTP are biased and noisy; best at start-up and on stable links \\
\key{Buffer-based} (BBA, BOLA) & Choose bitrate from buffer level only, e.g.\ BBA: \key{reservoir 5\,s}, cushion 10\,s (highest bitrate once the buffer exceeds 15\,s); BOLA uses \key{Lyapunov optimisation} and can abandon a chunk download & Robust in steady state, but ignores useful throughput information \\
\key{MPC} / robustMPC & Use throughput forecast + buffer to \key{maximise QoE over a horizon of 5 future chunks}; robustMPC divides the forecast by the max prediction error of the last 5 chunks & Needs an \key{accurate model} (throughput forecast); sensitive to prediction errors and to the horizon length \\
Tabular RL & Learn a table of values for every (state, action) & Table cannot scale; forced to assume e.g.\ \key{Markovian bandwidth} (only one past sample) \\
\bottomrule
\end{tabularx}
\end{center}
\textbf{Two failure examples of robustMPC (from the paper):}
\begin{itemize}
\item \textbf{Highly variable link (mean about 4.5\,Mb/s):} robustMPC's conservative estimate hovered near 2\,Mb/s, so it never used the high bitrates even as the buffer kept growing. The learned policy saw the high average and, once the buffer was large enough, switched to the top bitrate.
\item \textbf{HD-loving QoE metric:} the right strategy is to first build a big buffer at the lowest bitrate (300\,kb/s), then jump to HD (1850\,kb/s) and stay there. robustMPC kept a medium buffer and chose in-between bitrates because its 5-chunk horizon was too short to see the payoff. The learned policy did exactly the long-term plan (stayed at HD for nearly 80\,s).
\end{itemize}
\textbf{Lesson:} fixed heuristics or simplified models cannot (1) account for throughput fluctuations and prediction errors, or (2) choose the right planning horizon. Learning from the actual outcomes of decisions can.

\section{The QoE objective}
For a video of $N$ chunks, with $R_n$ the bitrate of chunk $n$, $q(\cdot)$ the perceived quality of a bitrate, and $T_n$ the rebuffering time caused by downloading chunk $n$:
\[ QoE = \underbrace{\sum_{n=1}^{N} q(R_n)}_{\text{quality}} \;-\; \mu \underbrace{\sum_{n=1}^{N} T_n}_{\text{rebuffering}} \;-\; \underbrace{\sum_{n=1}^{N-1} \big| q(R_{n+1}) - q(R_n) \big|}_{\text{smoothness}} \]
\begin{center}
\begin{tabularx}{\linewidth}{@{}lXl@{}}
\toprule
\textbf{Metric} & \textbf{Quality $q(R)$} & \textbf{Rebuffer penalty $\mu$} \\
\midrule
$QoE_{lin}$ (used by MPC) & $R$ (Mb/s) & 4.3 \\
$QoE_{log}$ (used by BOLA) & $\log(R/R_{min})$: diminishing returns at high bitrate & 2.66 \\
$QoE_{hd}$ & 0.3$\to$1, 0.75$\to$2, 1.2$\to$3, 1.85$\to$12, 2.85$\to$15, 4.3$\to$20 (HD gets a big jump) & 8 \\
\bottomrule
\end{tabularx}
\end{center}
\textbf{Worked example ($QoE_{lin}$):} four chunks at 1.2, 1.85, 1.85, 0.75\,Mb/s; chunk 3 causes 0.5\,s of rebuffering.
\begin{center}
\begin{tabular}{@{}lr@{}}
\toprule
Quality: $1.2 + 1.85 + 1.85 + 0.75$ & $5.65$ \\
Rebuffering: $4.3 \times 0.5$ & $-2.15$ \\
Smoothness: $|1.85-1.2| + 0 + |0.75-1.85|$ & $-1.75$ \\
\midrule
QoE & \key{1.75} \\
\bottomrule
\end{tabular}
\end{center}
Pensieve gives this per-chunk QoE to the learner as the \key{reward} $r_t$, so changing the metric changes the learned policy with no change to the algorithm.

\section{ABR as a reinforcement-learning problem}
\paperfig[0.6]{ps_fig2_rl}{RL for bitrate adaptation: the agent observes bandwidth, bitrate and buffer signals, the neural network outputs a bitrate, and the resulting QoE comes back as the reward.}
\begin{itemize}
\item At each step $t$ the \key{agent} observes a \key{state} $s_t$, takes an \key{action} $a_t$, the environment moves to $s_{t+1}$, and the agent gets a \key{reward} $r_t$.
\item Goal: maximise the expected cumulative \key{discounted} reward $\mathbb{E}\left[\sum_{t\ge0} \gamma^t r_t\right]$, $\gamma \in (0,1]$.
\item For ABR: one step = one chunk. State = player and network observations; action = bitrate of the next chunk; reward = that chunk's QoE.
\item The policy is \key{not hand-written}; it is a neural network trained from experience. It can learn, for a given network, how much buffer is needed against rebuffering, how much to trust throughput vs.\ buffer signals, and how far ahead to plan.
\end{itemize}

\section{Training environment: a chunk-level simulator}
Training on real players is too slow (a browser must download every chunk before the model can update). Pensieve uses a simple simulator that still captures the player's dynamics.

\textbf{Per chunk download, the simulator:}
\begin{enumerate}
\item Computes the \key{download time} from the chunk's size (its bitrate) and the \key{throughput trace}.
\item \key{Drains} the playback buffer by that download time (video keeps playing while downloading).
\item Records \key{rebuffering} if the download time exceeds the buffer at the start of the download.
\item \key{Adds} the chunk's playback duration to the buffer.
\item If the buffer cannot fit another chunk, \key{waits 500\,ms} and retries (the default DASH retry rate).
\item Reports to the agent: buffer level, rebuffering time, download time, next chunk's sizes at all bitrates, chunks remaining.
\end{enumerate}
Result: \key{100 hours of video in 10 minutes} of training.

\textbf{Worked example} (4\,s chunks):
\begin{center}
\begin{tabular}{@{}lccccc@{}}
\toprule
\textbf{Chunk} & \textbf{Size} & \textbf{Throughput} & \textbf{Download} & \textbf{Rebuffer} & \textbf{Buffer after} \\
\midrule
start & & & & & 6.00\,s \\
1.85\,Mb/s & 7.4\,Mb & 2.5\,Mb/s & 2.96\,s & 0 & $6 - 2.96 + 4 = 7.04$\,s \\
4.3\,Mb/s & 17.2\,Mb & 2.0\,Mb/s & 8.60\,s & $8.6 - 7.04 = 1.56$\,s & $0 + 4 = 4.00$\,s \\
\bottomrule
\end{tabular}
\end{center}

\subsection{The transport-layer catch: slow-start restart}
\paperfig[0.48]{ps_fig4_ssr}{Per-chunk throughput on a 6\,Mb/s link. With TCP slow-start restart (a), low-bitrate (small) chunks get much less throughput; with it disabled (b), every bitrate gets about 6\,Mb/s.}
\begin{itemize}
\item A player does not always request the next chunk immediately (e.g.\ the buffer is full). After such idle time TCP goes back to \key{slow start} (\key{slow-start restart}), so small chunks never reach full speed.
\item Then throughput depends on the \emph{ABR algorithm itself} (algorithms that fill the buffer fast pause more and get more slow starts), which a trace-driven simulator cannot model.
\item Fix: \key{disable slow-start restart on the video server} (Linux: \texttt{net.ipv4.tcp\_slow\_start\_after\_idle}). Then each chunk uses the full trace throughput, as the simulator assumes. This can make traffic burstier; standards work on pacing the first burst after idle addresses that.
\item No simulator is perfect, but Pensieve learns good policies from an imperfect one \key{as long as training covers a wide variety of network conditions} (it generalises well).
\end{itemize}

\section{The learning algorithm: A3C (actor-critic)}
\paperfig[0.62]{ps_fig5_a3c}{Pensieve's two networks. The state goes through 1D-CNNs (past throughputs, download times, next chunk sizes) and fully connected layers. The actor outputs a probability for each bitrate; the critic outputs the value of the state.}

\subsection{State (input), after chunk $t$ is downloaded}
\[ s_t = (\vec{x}_t,\ \vec{\tau}_t,\ \vec{n}_t,\ b_t,\ c_t,\ l_t) \]
\begin{center}
\begin{tabular}{@{}ll@{}}
\toprule
$\vec{x}_t$ & throughput measured for the past $k$ chunks \\
$\vec{\tau}_t$ & download times of the past $k$ chunks (the time span of each throughput sample) \\
$\vec{n}_t$ & sizes of the next chunk at each of the $m$ bitrates \\
$b_t$ & current buffer level \\
$c_t$ & number of chunks left in the video \\
$l_t$ & bitrate of the last chunk \\
\bottomrule
\end{tabular}
\end{center}

\subsection{Policy (actor)}
The \key{policy} $\pi_\theta(s_t, a_t) \in [0,1]$ is the probability of choosing bitrate $a_t$ in state $s_t$. States are continuous (throughput, buffer), so a table is impossible; a \key{neural network} with parameters $\theta$ represents the policy. NNs need \key{no hand-crafted features}: they take raw signals.

\subsection{Training the actor: policy gradient}
The gradient of the expected discounted reward is
\[ \nabla_\theta\, \mathbb{E}_{\pi_\theta}\Big[\sum_{t} \gamma^t r_t\Big] = \mathbb{E}_{\pi_\theta}\big[\nabla_\theta \log \pi_\theta(s,a)\, A^{\pi_\theta}(s,a)\big] \]
where the \key{advantage} $A(s,a)$ = how much better action $a$ is than the policy's average action in state $s$. The update after a sampled trajectory:
\[ \theta \leftarrow \theta + \alpha \sum_t \nabla_\theta \log \pi_\theta(s_t,a_t)\, A(s_t,a_t) \;+\; \beta\, \nabla_\theta H\big(\pi_\theta(\cdot \mid s_t)\big) \]
\begin{itemize}
\item $\nabla_\theta \log \pi_\theta(s_t,a_t)$ is the direction that makes action $a_t$ \key{more likely} in $s_t$; the step size is the advantage. Net effect: \key{reinforce actions that turned out better than average}, discourage the others.
\item The \key{entropy} term $H$ pushes towards more random policies to keep \key{exploring}. $\beta$ starts large and \key{decays from 1 to 0.1 over $10^5$ iterations}, shifting from exploration to exploitation.
\end{itemize}

\subsection{Training the critic: temporal difference}
The \key{critic} learns the \key{value} $V(s)$ (expected total reward from $s$ under the policy):
\[ \theta_v \leftarrow \theta_v - \alpha' \sum_t \nabla_{\theta_v} \big(r_t + \gamma V(s_{t+1};\theta_v) - V(s_t;\theta_v)\big)^2 \]
and the advantage is estimated as $A(s_t,a_t) \approx r_t + \gamma V(s_{t+1}) - V(s_t)$.

\textbf{Worked example:} $r_t = 2$, $\gamma = 0.99$, $V(s_{t+1}) = 10$, $V(s_t) = 11$. Then $A = 2 + 9.9 - 11 = 0.9 > 0$: the chosen bitrate did better than expected, so its probability is pushed up. Had the chunk caused rebuffering ($r_t = -3$), $A = -4.1$ and the probability would be pushed down.

\textbf{Only the actor is used after training}; the critic exists just to train it.

\subsection{Parallel, asynchronous training}
\begin{itemize}
\item \key{16 agents} run in parallel, each on different network traces.
\item Each sends its \{state, action, reward\} tuples to a \key{central agent}, which computes the actor-critic gradient, updates the model, and pushes the new model back to that agent.
\item Fully \key{asynchronous} (no locking between agents).
\item Why A3C: state of the art, and its asynchronous structure fits a future where \key{many real clients} send experience concurrently.
\end{itemize}

\section{One model for many videos}
\paperfig[0.48]{ps_fig6_mask}{Canonical input and output: chunk sizes are placed at the index of the closest canonical bitrate (unsupported ones are 0), and a 0-1 mask removes unsupported bitrates from the softmax.}
Videos differ in how many bitrates they offer and in their chunk sizes, but an NN has fixed-size inputs and outputs. Training one model per video type does not scale. Two changes:
\begin{enumerate}
\item \key{Canonical input}: fix a maximum number of bitrate levels (e.g.\ \key{13} covers all DASH reference videos). Put each chunk size at the index of the \key{closest bitrate}; \key{zero} the rest.
\item \key{Masked softmax output}: with mask $m_i \in \{0,1\}$ (1 if the video supports bitrate $i$) and network outputs $z_i$,
\[ p_i = \frac{m_i\, e^{z_i}}{\sum_j m_j\, e^{z_j}} \]
The mask depends only on the video, not on $\theta$, so $p_i$ stays differentiable and training is unchanged.
\end{enumerate}
\textbf{Worked example:} mask $(1,1,0,1,0)$, outputs $z = (1.0,\ 2.0,\ 3.0,\ 0.5,\ 4.0)$.
\begin{center}
\begin{tabular}{@{}lccccc@{}}
\toprule
Index $i$ & 1 & 2 & 3 & 4 & 5 \\
\midrule
$m_i e^{z_i}$ & 2.718 & 7.389 & 0 & 1.649 & 0 \\
$p_i$ (sum = 11.756) & 0.231 & \key{0.629} & 0 & 0.140 & 0 \\
\bottomrule
\end{tabular}
\end{center}
Without the mask, index 5 ($z=4$) would get most of the probability although this video does not offer that bitrate.

\section{Implementation and deployment}
\begin{center}
\begin{tabularx}{\linewidth}{@{}lX@{}}
\toprule
\textbf{Item} & \textbf{Setting} \\
\midrule
History & $k = 8$ past throughput samples \\
Throughput input & 1D-CNN, \key{128 filters of size 4}, stride 1 \\
Chunk sizes input & another 1D-CNN of the same shape \\
Merge & hidden layer of \key{128 neurons}, then softmax over bitrates (actor) \\
Critic & same structure, final output one \key{linear} neuron \\
Discount $\gamma$ & \key{0.99} (an action is influenced by about 100 future steps) \\
Learning rates & actor $10^{-4}$, critic $10^{-3}$ \\
Entropy weight $\beta$ & 1 $\to$ 0.1 over $10^5$ iterations \\
Software & TensorFlow (TFLearn API) \\
\bottomrule
\end{tabularx}
\end{center}
Performance was found to be insensitive to these hyperparameters over a wide range, so no sophisticated tuning was used.

\textbf{Deployment: an ABR server.}
\begin{enumerate}
\item The trained actor runs on a standalone \key{ABR server} (Python \texttt{BaseHTTPServer}).
\item Each client chunk request \key{carries the observations} (last chunk's throughput and download time, buffer, video being streamed, \ldots).
\item The server runs the actor NN and replies with the \key{bitrate for the next chunk}.
\item The client fetches that chunk from the CDN.
\end{enumerate}
\begin{itemize}
\item The server stays \key{stateless}: everything it needs is in the request.
\item Clients with little compute (TVs, phones) need not run a neural network; the observations still come from the client, so it reacts as quickly as a client-side ABR. Adding 100\,ms RTT to reach the server cost only about 3.5\% QoE.
\item The model could also run inside players (compressed NNs, or JavaScript).
\end{itemize}

\section{Discussion and limits}
\begin{itemize}
\item \key{Periodic or online training}: models could be retrained as new data arrives. Training \emph{on the client} is harder: it costs client compute and needs algorithms that learn from little data quickly.
\item \key{Why not tabular RL}: a table needs few, quantized states, which forces assumptions like Markovian bandwidth (only the last throughput sample). One sample is not enough to infer the distribution of future bandwidth; a neural network can use the whole history.
\item \key{Optimality}: there is still a gap to the \emph{offline} optimum, but that optimum knows future bandwidth exactly. Against the optimal \emph{online} policy (computable by dynamic programming for a known Markov model of download times), Pensieve comes much closer.
\item Training data: throughput traces (broadband and 3G/HSDPA); traces where the choice is trivial (always max bitrate, or nothing feasible) were excluded. A model trained only on \key{synthetic traces} still generalised to real networks.
\end{itemize}

\begin{keybox}[Exam one-liner]
Pensieve replaces hand-tuned ABR rules with a policy learned by deep RL: an actor-critic (A3C) network takes past throughputs and download times, next-chunk sizes, buffer, chunks left and last bitrate, outputs a probability per bitrate, and is trained in a chunk-level simulator to maximise per-chunk QoE (quality $-$ $\mu\cdot$rebuffering $-$ switches); masking gives one model for many bitrate ladders, and an ABR server serves decisions to clients.
\end{keybox}

\section{Revision checklist}
\begin{itemize}
\item \textbf{DASH:} chunks of a few seconds, several aligned bitrates, manifest $\to$ CDN, ABR picks each chunk's bitrate, playback only after full download.
\item \textbf{Challenges:} variable throughput, conflicting QoE goals, cascading effects, coarse choices.
\item \textbf{Baselines:} rate-based (harmonic mean of 5), buffer-based (BBA 5\,s reservoir; BOLA Lyapunov), MPC (5-chunk horizon, needs forecasts), tabular RL (Markov assumption).
\item \textbf{QoE:} $\sum q(R_n) - \mu\sum T_n - \sum|q(R_{n+1})-q(R_n)|$; lin ($\mu$ 4.3), log ($\mu$ 2.66), hd ($\mu$ 8).
\item \textbf{Simulator:} download time from size/trace, drain buffer, count rebuffer, add chunk, 500\,ms wait if full; 100\,h in 10\,min; disable TCP slow-start restart.
\item \textbf{State:} $(\vec x, \vec\tau, \vec n, b, c, l)$, $k=8$.
\item \textbf{A3C:} actor $\pi_\theta$ updated by $\nabla\log\pi\cdot A$ + entropy (1$\to$0.1); critic by TD; $A \approx r + \gamma V(s') - V(s)$; only actor used at run time; 16 async agents.
\item \textbf{NN:} 1D-CNN 128 filters size 4; 128-neuron hidden layer; softmax; $\gamma=0.99$; LR $10^{-4}$/$10^{-3}$.
\item \textbf{Many videos:} 13 canonical levels, nearest-bitrate mapping, masked softmax $p_i = m_i e^{z_i}/\sum m_j e^{z_j}$.
\item \textbf{Deployment:} stateless ABR server; observations in each request; reply = next bitrate.
\end{itemize}
